\clearpage
\fancyhead[L]{\small BODE. ASTRONOMISCHES JAHRBUCH.}
\hypertarget{native-581-lower}{}
\begin{center}\small Göttingische gelehrte Anzeigen. Stück 28. Seiten 265..277. 18. Februar 1815.\\[3pt]\rule{.72\linewidth}{.3pt}\end{center}

\emph{Astronomisches Jahrbuch für das Jahr} 1817 \emph{nebst einer Sammlung der neuesten Abhandlungen, Beobachtungen und Nachrichten aus den astronomischen Wissenschaften. Berechnet und herausgegeben von} J. E. \textsc{Bode}, \emph{königlichem Astronomen und Mitglied der Akademie.} Beim Verfasser und in Kommission bei J. E. \textsc{Hitzig}. Berlin 1814. 260 Seiten im Oktavformat, mit einer Kupfertafel.

Im Jahr 1817 gibt es nur zwei Sonnenfinsternisse, die beide in Europa unsichtbar sind. Unter den Erscheinungen, die die besondere Aufmerksamkeit der Astronomen verdienen, heben wir eine sehr enge Begegnung der Venus mit Regulus hervor. Nach der Berechnung von Herrn \textsc{Bode} wird sie am 28. September um $14^{\mathrm{u}}$ Berliner Zeit stattfinden und sich daher im östlichen Teil Europas sehr gut beobachten lassen. Nach Herrn \textsc{Bode}s Rechnung zieht die Venus nur eine Bogenminute südlich am Stern vorbei, sodass es beinahe zu einer wirklichen Bedeckung kommt. Solche Erscheinungen verdienen es, im Voraus so genau wie möglich berechnet zu werden. \emph{Wirkliche} Bedeckungen der helleren Sterne durch die Venus, besonders wenn sie sich näher an der unteren Konjunktion mit der Sonne befindet, würden nach unserer Überzeugung ebenso sehr weite Reisen rechtfertigen wie die Durchgänge der Venus vor der Sonne. Bemerkenswert ist auch eine ebenso enge, aber in Europa unsichtbare Begegnung des Mars mit Saturn am 18. April um $7^{\mathrm{u}}$, die man der Sternwarte in Botanybay zur Beobachtung empfehlen könnte.

Die \emph{Zusätze} zum vorliegenden Band des Jahrbuchs beginnen mit Ephemeriden für Pallas von Herrn \textsc{Nicolai}, für Juno von Herrn \textsc{Möbius} und für Vesta von Herrn \textsc{Gerling}. Es folgen Beobachtungen und parabolische Elemente des zweiten Kometen von 1813 von Herrn Dr.\ \textsc{Olbers}. --- Astronomische Beobachtungen an der Prager Sternwarte im Jahr 1813 von den Herren \textsc{David} und \textsc{Bittner}. Das hier empfohlene Verfahren, die Azimute ohne Heranziehung der Zeit aus einander entsprechenden horizontalen Abständen von einem Stern bei gleichen Höhen zu bestimmen, ist an sich sehr gut. Bei einem Instrument, dessen horizontaler und vertikaler Kreis beide hinreichend groß und fein unterteilt sind, wäre es auch für sehr genaue Bestimmungen geeignet. Verwendet man aber einen Theodoliten, dessen Höhenkreis nur in Bogenminuten unterteilt ist, oder einen Kreis mit einer solchen Teilung des Azimutkreises, so ist natürlich nur eine geringere Genauigkeit zu erwarten. Beim ersten Instrument ließe sich dieser Nachteil zwar vermeiden, wenn man das Fernrohr unverrückt auf der beobachteten Höhe stehen ließe. Dann fiele aber der Vorteil der Vervielfachung weg. Beim Kreis würden sich außerdem kleine Fehler in der Justierung der Querlibelle und in der Ausrichtung der Sichtlinie parallel zur Ebene des Instruments in voller Stärke auf das Ergebnis auswirken. Sie sind kaum ganz zu vermeiden, bleiben aber bei der Verwendung des Instruments für seinen eigentlichen Zweck unschädlich. Diese Fehler lassen sich jedoch aufheben, wenn man am folgenden Tag dieselben Beobachtungen bei entgegengesetzter Lage des Kreises wiederholt, da sie dann \SourceLeaf{582}auch in entgegengesetzter Richtung wirken müssen. --- Weitere Bemerkungen und Zusätze zu den Beobachtungen über $\eta$ Antinous im astronomischen Jahrbuch für 1816 von Herrn Prof.\ \textsc{Wurm}. Dieser geschickte Astronom zog zehn ältere Beobachtungen von \textsc{Pigott} heran, die Herrn \textsc{Wurm} bisher entgangen waren, sowie einige weitere eigene. Mit der Methode der kleinsten Quadrate bestimmte er nun die Periode zu 7,17604 Tagen und die Epoche zu 1800 Januar 4,504. Zugleich werden auf Grundlage dieser Bestimmung neue Tafeln mitgeteilt. --- Mehrere wertvolle astronomische Beiträge von Herrn Prof.\ \textsc{Littrow} in Kasan betreffen die Entwicklung der Mittelpunktsgleichung in eine Reihe und die Theorie der Epizykel. Mit Epizykeln lassen sich, wie der Verfasser sehr gut zeigt, elliptische Bewegungen nicht zugleich in der Länge und im Radiusvektor darstellen. Es geht ferner um ein Verfahren, die Kreisbahn eines Planeten aus zwei Beobachtungen zu bestimmen. Dieses ist allgemeiner und der Form nach einfacher als das vom selben Verfasser im Jahrbuch 1816 entwickelte. Bei mäßigen Bewegungen würde es aber wegen des eingeführten Kosinus dieser Bewegungen in der Anwendung eine geringere Genauigkeit ergeben, obwohl die ganze Aufgabe gegenwärtig freilich nur theoretisches Interesse hat. Dieselbe Bemerkung gilt für die Bestimmung einer geradlinigen Bahn aus vier beobachteten Längen (dieser Aufsatz ist allein durch Druckfehler völlig entstellt). Bei derselben Aufgabe für drei vollständige Beobachtungen wendet der Verfasser die aus der \emph{Theoria motus corporum coelestium} entnommene Lösung auf drei Beobachtungen der Vesta an (die Voraussetzung ist an sich nicht einmal als Näherung zulässig, wie im angeführten Werk, Art. 131, gezeigt wird. Herr \textsc{Littrow} schlägt sie zwar nur vor, um einen neu entdeckten Himmelskörper nach einer kurzen Unterbrechung leichter wiederzufinden. Dafür werden jedoch schon die gewöhnlichen Interpolationsverfahren ausreichen). Herr \textsc{Littrow} gibt ferner eine indirekte Lösung für die Aufgabe, aus den Höhen zweier Sterne die Zeit und die Polhöhe zu finden. Sie ist an sich zweckmäßig und wird noch bequemer, wenn man logarithmische Differenzen verwendet. Außerdem ist anzumerken, dass man, sooft die Stundenwinkel $\alpha$, $\alpha'$ kleiner als $90^{\circ}$ sind, $\frac{1}{2}\alpha$ und $\frac{1}{2}\alpha'$ besser nach den bekannten Formeln aus ihren Sinuswerten bestimmt. Bei der Berechnung von Zirkummeridianhöhen würden wir das gewöhnliche Verfahren dem neuen vorziehen. Der Gewinn an Genauigkeit ist zu gering, da der Fehler von derselben Ordnung bleibt wie beim gewöhnlichen Verfahren, wenn man sich auf das erste Glied beschränkt. Ebenso ziehen wir bei der letzten Aufgabe zum Erdsphäroid die gleichwertigen endlichen Ausdrücke den hier gegebenen Reihen vor. --- Es folgen astronomische Beobachtungen an der Sternwarte in Wien von den Herren \textsc{Triesnecker} und \textsc{Bürg} und an der Sternwarte in Cremsmünster von Herrn \textsc{Derfflinger}. --- Beobachtungen zur Bestimmung der geografischen Lage von Port Jackson, von Herrn J. \textsc{Oltmanns} aus \textsc{Malaspinas} Papieren berechnet. Hier wird uns Hoffnung auf eine vollständige Bearbeitung sämtlicher Beobachtungen gemacht. --- Die Ideen zur Störungsrechnung nach \textsc{Keppler} von Herrn Prof.\ J. W. \textsc{Pfaff} in Nürnberg scheinen uns doch zu oberflächlich, um irgendein Ergebnis zu liefern. Die hier aufgestellte Form der Störungsgleichungen für Knoten und Neigung (oder vielmehr für ihre differenziellen Änderungen) scheint uns ganz unfruchtbar. Insbesondere folgen daraus keine Beziehungen der Art, die Herr \textsc{Pfaff} annimmt, da zum Beispiel aus den Gliedern $Fxx'$ und $Gyy'$ Gleichungen mit derselben Periode hervorgehen. --- Beobachtungen an der Wilnaer Sternwarte von Herrn Prof.\ \textsc{Sniadecky}. --- Die Bemerkungen von Herrn Prof.\ \textsc{Bürg} über geografische Ortsbestimmungen in Ungarn, Österreich und Bayern enthalten eine Verteidigung gegen einen Aufsatz in der monatlichen Correspondenz. Dieser hatte insbesondere die in Wien vorgenommene Azimutbestimmung infrage zu stellen versucht. --- Über den Kometen von 1558 von Herrn Dr.\ \textsc{Olbers}. Bisher kannte man von diesem Kometen nur drei Beobachtungen des Landgrafen \textsc{Wilhelm} vom 20., 21. und 23. September und eine von \textsc{Cornelius Gemma} vom 20. September. Diese letzte wich stark von \SourceLeaf{583}der gleichzeitigen Beobachtung des Landgrafen ab. Durch eine sehr glückliche Verbesserung des Textes bei \textsc{Gemma}, indem er distabat eo die statt distabat 20 die liest, löst Herr Dr.\ \textsc{Olbers} diesen Widerspruch auf die befriedigendste Weise. \textsc{Gemma}s Beobachtung gilt nun für den 17. August und ermöglicht wenigstens eine annähernde Bestimmung der Kometenbahn, die hier mitgeteilt wird. --- Einige physikalisch-astronomische Beobachtungen von Saturn, Mars, Mond, Venus und Sonne von Herrn Dr.\ \textsc{Gruithusen} in München geben bemerkenswerte Belege für die Vorzüglichkeit der Fernrohre von \textsc{Fraunhofer}. --- Ebenfalls von Herrn Prof.\ \textsc{Oltmanns} ein Beitrag zur Bestimmung der geografischen Länge von \textsc{Quito} aus verschiedenen Beobachtungen der französischen Akademiemitglieder und die Berechnung der kürzlich zur Sprache gebrachten totalen Sonnenfinsternis in Mirabeau in der Provence am 3. Juni 1239. Herr \textsc{Oltmanns} wendet sie mit Erfolg auf die Bewegung des Mondknotens an. Die in den Tafeln von Herrn \textsc{Oltmanns} angesetzte Bewegung muss demnach um $1'40''$ auf 100 Jahre verringert werden. Unter Heranziehung zweier anderer ähnlicher Beobachtungen findet Herr \textsc{Oltmanns} für die Verringerung der Bewegung in hundert Jahren im Mittel $1'14''$. --- Die Tafeln für die scheinbaren Örter des Polarsterns von Herrn Prof.\ \textsc{Bessel} sind für die beobachtenden Astronomen ein außerordentlich wertvolles Geschenk. Es ist schade, dass der Verfasser dabei die neueren Beobachtungen der Rektaszension dieses Sterns durch verschiedene Astronomen noch nicht verwenden konnte. Sie weisen übereinstimmend nach, dass eine Vergrößerung nötig ist (gegenwärtig um etwa 4 Zeitsekunden). Wir hätten uns gewünscht, dass auch die reinen mittleren Positionen für jedes Jahr beigefügt wären, damit man die Bestimmungen anderer Astronomen leichter vergleichen kann. Bemerkungen von Herrn Prof.\ \textsc{Bürg} über die Revision seiner früheren Mondberechnungen. Um seine Mondtafeln nochmals zu verbessern, hat Herr \textsc{Bürg} begonnen, die Beobachtungen in Greenwich seit 1765 mit ihnen zu vergleichen. Wie auch immer die Ergebnisse dieser mühsamen Arbeit ausfallen, sie werden gewiss interessant und lehrreich sein. --- Beobachtungen der Juno und neue Elemente ihrer Bahn, Berechnung der nächsten Opposition der Pallas und andere astronomische Nachrichten von Herrn Prof.\ \textsc{Gauss}. --- Einiges über die Aussicht auf neue Entdeckungen am Himmel durch Fernrohre. Hier werden die Schwierigkeiten dargelegt, die einer weiteren Verbesserung der optischen Instrumente Grenzen setzen. Wir gestehen jedoch, dass wir dadurch noch nicht davon überzeugt sind, dass diese Hindernisse völlig unüberwindlich sind. Inwieweit die hier hervorgehobenen Mängel den Nutzen der Vorschläge mindern oder sogar aufheben, das Fernrohr durch ein Uhrwerk der Bewegung der Sterne folgen zu lassen, müssen wir offenlassen. Auch verstehen wir nicht ganz, wie man auf S.\,218 behaupten kann, dass das ständige Hin- und Herfliegen des Bildes eine ruhige Wahrnehmung \emph{während der wenigen Sekunden seines Durchgangs} unmöglich mache, da die Vorrichtung doch gerade bewirkt, dass das Bild lange im Gesichtsfeld bleibt. Sollte die Erfahrung jedoch das Gewicht dieses Einwands bestätigen, so könnte die Mechanik, wie uns scheint, wohl Mittel finden, auch ohne Zahn und Getriebe. dem Fernrohr eine äußerst sanfte und gleichförmige Bewegung zu geben. Das Haupthindernis für eine weitere Steigerung der Leistung optischer Instrumente dürfte in den atmosphärischen Schwankungen liegen, die umso deutlicher werden, je stärker jene vergrößern. --- Über zwei veränderliche Sterne im Herkules von Herrn Dr.\ \textsc{Koch} in Danzig. Leider ist auch die Danziger Sternwarte, aus der so manche nützliche Beobachtung hervorgegangen ist, ein Opfer der letzten Belagerung geworden. Der in der Anmerkung von Herrn \textsc{Bode} mit $m$ bezeichnete Stern, der noch in keinem Verzeichnis steht, ist übrigens schon von unserem Herrn Prof.\ \textsc{Harding} bestimmt worden und findet sich auf Blatt XVI seiner Himmelskarten. --- Die mittleren Positionen von 28 der bedeutendsten Sterne der Plejaden nach dem neuen Sternkatalog von \textsc{Piazzi}, aus demselben Katalog die mittleren Positionen der 36 \textsc{Maskelyne}-Fundamentalsterne und schließlich die des Polarsterns. Besonders bemerkens\SourceLeaf{584}wert ist bei diesem die aus den Beobachtungen von Herrn \textsc{Cacciatore} folgende Parallaxe in Rektaszension von $2^{\mathrm{s}}885$ in Zeit. Sie scheint uns allerdings zu groß und lässt sich mit den zahlreichen Beobachtungen, die Herr \textsc{von Lindenau} an der Seeberger Sternwarte angestellt hat, nicht vereinbaren. Aus den von \textsc{Oriani} beobachteten Zenitdistanzen hatten wir nur eine halb so große Parallaxe abgeleitet (siehe das 10. Stück dieser Zeitschrift [S. 574 dieses Bandes]), die freilich ebenfalls keine absolute Zuverlässigkeit hat. Die jährliche Eigenbewegung des Polarsterns in Rektaszension findet \textsc{Piazzi} aus dem Vergleich

\begin{center}
\small\begin{tabular}{@{}l@{\hspace{8pt}}c@{\hspace{8pt}}r@{}}
mit Hevel & $.\;.\;.\;.\;.\;.\;.\;.\;.\;.\;.\;.\;.$ & $+6^{\mathrm{s}}82$\\
Flamstead & $.\;.\;.\;.\;.\;.\;.\;.\;.\;.\;.\;.$ & $+9.03$\\
La Caille & $.\;.\;.\;.\;.\;.\;.\;.\;.\;.\;.\;.$ & $+3.96$\\
Bradley & $.\;.\;.\;.\;.\;.\;.\;.\;.\;.\;.\;.$ & $+1.62$\\
\end{tabular}
\end{center}
\noindent in Zeit. Herr \textsc{von Zach} hatte aus dem Vergleich der Bestimmungen von \textsc{La Caille} und \textsc{Mathieu} $+3^{\mathrm{s}}177$ gefunden. Es wäre sehr zu wünschen, dass dieser wichtige Gegenstand einmal sorgfältig untersucht würde. Dabei müssten aber vorher alle Originalbeobachtungen auf einheitliche Weise reduziert werden. --- Astronomische Beobachtungen an der Königl Sternwarte in Berlin im Jahr 1813. --- Beobachtungen von Herrn Prof.\ \textsc{Struve} an der Sternwarte in Dorpat in den Jahren 1812, 1813 und 1814. --- Weitere verschiedene Beobachtungen in Cremsmünster. --- Ferner Nachrichten über den Doppelstern 61 Cygni aus \textsc{Piazzi}s Libro sesto del reale osservatorio (wo die bemerkenswerte Eigenbewegung dieses Doppelsterns zuerst mitgeteilt wurde) und aus dem neuen großen Sternkatalog dieses Astronomen. --- Nachweis, dass keiner von acht am Himmel vermissten Fixsternen Ceres, Pallas, Juno oder Vesta war. Wir bemerken hierzu, dass alle diese Sterne bis auf den in der nordlichen Krone bereits im neunten Band der monatlichen Correspondenz (S. 153. 151. 241. 155. 243. 241. 246) erörtert und insbesondere alle hier angeführten Sterne von \textsc{Mayer} aus den Originalpapieren bereits befriedigend erklärt worden sind. ---

Unter den vermischten Bemerkungen und Nachrichten, die wie gewöhnlich am Ende des Jahrbuchs angehängt sind, heben wir hier nur eine hervor. Es ist eine Stelle aus einem Brief von Herrn Dr.\ \textsc{Reche} in Mühlheim am Rhein vom 6. Mai 1814. Dem Brief zufolge hat Herr Prof.\ \textsc{Kramp} während der Blockade von Straßburg das Problem, die Bahn eines Kometen aus drei Beobachtungen zu bestimmen, auf eine neue Weise gelöst. Er hat gefunden, dass alle von ihm berechneten Kometenbahnen, die man bisher für Parabeln gehalten hatte, \emph{ganz gewiss und bestimmt} Hyperbeln sind. Eine solche Ankündigung muss sehr auffallen. Da die Arbeit von Herrn \textsc{Kramp} seitdem in der Juliausgabe der mathematischen Zeitschrift erschienen ist, die seit einigen Jahren in Nimes unter dem Titel Annales de mathématiques pures et appliquées herauskommt, ist dies der Ort, die vorangehende Besprechung mit einer Beurteilung dieser Abhandlung zu verbinden.

Hier kommt es auf zwei Punkte an. Der eine ist der Wert des Verfahrens an sich. Der andere ist die Art seiner Anwendung auf tatsächlich beobachtete Kometen und die Zuverlässigkeit der daraus hervorgehenden Ergebnisse. Wir wollen zuerst einiges über die letzteren Punkte bemerken. In dieser Abhandlung ist nur von der Anwendung auf einen einzigen Kometen die Rede, nämlich auf den von 1781. Aus den Beobachtungen vom 14, 19. und 25. November leitet Herr \textsc{Kramp} eine Hyperbel mit der Exzentrizität $= 4,580652$ und der Periheldistanz $1,048364$ ab. Wer auch nur einige Erfahrung in der Berechnung von Kometenbahnen hat, dem würde schon der kurze Zeitabstand genügen, um die Zuverlässigkeit dieses Ergebnisses zu beurteilen. Wir müssen aber hinzufügen, dass die gefundenen Elemente nicht einmal mit den drei zugrunde liegenden Beobachtungen selbst verglichen werden, geschweige denn an anderen Beobachtungen geprüft werden. Wer eine solche Prüfung selbst vornehmen wollte, dem wird sie zumindest dadurch sehr erschwert, dass Herr \textsc{Kramp} die Berechnung der Elemente nicht einmal zu \SourceLeaf{585}Ende geführt hat, sodass die Länge des Perihels und der Zeitpunkt des Durchgangs fehlen. Wenn man schließlich noch bedenkt, dass die Erdbahn in der ganzen Rechnung als Kreis behandelt wird, wird man es wohl für überflüssig halten, eigens zu untersuchen, ob vielleicht auch in der numerischen Rechnung selbst Fehler gemacht worden sind. Dafür scheint der \emph{bedeutende} Unterschied zwischen dem Wert der Größe
\[
P\!\cdot\!Q'' - P''Q = 0.4614411
\]
und der Summe der beiden Größen
\[
\begin{aligned}
PQ' - P'Q &= 0.2133547\\
P'Q'' - P''Q' &= 0.2495918
\end{aligned}
\]
zu sprechen. Diese Summe müsste sich eigentlich, wie aus dem Verfahren folgt, \emph{genau} gleich der ersten Größe ergeben. (Herr \textsc{Kramp} urteilt darüber anders und nennt diesen beträchtlichen Unterschied eine egalité presque rigoureuse, aber wozu mit sieben Dezimalstellen rechnen, wenn man nicht einmal für die dritte einstehen könnte?) Wir sagen \emph{scheint}, weil dieser Unterschied vielleicht eine Folge der Weglassungen auf S. 13 sein könnte.

Wir wenden uns nun dem Verfahren selbst zu. Es wurde schon darauf hingewiesen, dass Herr \textsc{Kramp} (tatsächlich ganz ohne Not) die Exzentrizität der Erdbahn völlig vernachlässigt. Das Wesentliche des Verfahrens lässt sich mit wenigen Worten angeben. Herr \textsc{Kramp} bezeichnet die halbe große Achse der Bahn des Himmelskörpers mit $b$. Die Koordinaten seiner drei Positionen in der Bahn, bezogen auf die Knotenlinie als Abszissenachse, bezeichnet er mit $bP$, $bP'$, $bP''$ sowie $bQ$, $bQ'$, $bQ''$. Diese sechs Größen lassen sich bekanntlich leicht aus den Beobachtungsdaten, der Bahnneigung $\beta$ und der Länge des aufsteigenden Knotens $\delta$ bestimmen. Andererseits kombiniert Herr \textsc{Kramp} die Formeln, die sich auf die Gesetze von \textsc{Kepler} beziehen, und setzt die Differenzen der exzentrischen Anomalien ihren Sinuswerten gleich. So findet er, dass die drei Größen
\[
\begin{gathered}
PQ' - P'Q\\
P'Q'' - P''Q'\\
PQ'' - P''Q
\end{gathered}
\]
der Reihe nach proportional zu den Zeitabständen $t$, $t'$, $t+t'$ sind, sodass sich die Gleichung
\[
0 = PQ' - P'Q + P'Q'' - P''Q' + P''Q - PQ''
\]
ergibt, die Herr \textsc{Kramp} eine équation essentielle et très remarquable par sa simplicité nennt. --- Setzt nun Herr \textsc{Kramp} in den beiden Gleichungen
\[
\begin{aligned}
(t+t')(PQ'-P'Q) &= t(PQ''-P''Q)\\
(t+t')(P'Q''-P''Q') &= t'(PQ''-P''Q)
\end{aligned}
\]
für $P$, $P'$, $P''$, $Q$, $Q'$, $Q''$ ihre Werte in Abhängigkeit von $\beta$ und $\delta$ ein (wobei $b$ offenbar herausfällt), so erhält er zwei Gleichungen zwischen $\beta$ und $\delta$. Aus diesen leitet er unter Weglassung einiger Glieder durch Elimination von $\beta$ schließlich eine kubische Gleichung für $\operatorname{tang}\delta$ ab. Aus dieser bestimmt man $\delta$, dann aus einer der beiden eben genannten Gleichungen (oder aus beiden) $\beta$, dann die Größen $bP$, $bP'$, $bP''$, $bQ$, $bQ'$, $bQ''$. Aus vier dieser Größen bestimmt man schließlich die übrigen Elemente des Kegelschnitts nach Verfahren, die wir der Kürze halber hier übergehen.

Es wird sehr leicht sein, die völlige Unzulässigkeit dieses Verfahrens zu beweisen. An die Stelle der Bewegungsgesetze tritt hier, wie wir sehen, die Proportionalität der mehrmals erwähnten drei Größen zu den Zeitabständen. Die ganze Lösung beruht allein auf dieser Proportionalität oder den beiden ihr gleichwertigen Gleichun\SourceLeaf{586}gen und auf dem Satz, dass die drei Positionen in einer Ebene liegen. Die drei Größen
\[
bb(PQ'-P'Q), \quad bb(P'Q''-P''Q'), \quad bb(PQ''-P''Q)
\]
sind, wie man sofort bemerken wird, genau die drei in der Theoria motus Corporum Coelestium, Art. 112, mit $n''$, $n$, $n'$ bezeichneten Größen. Dort wird aber in Art. 114 bewiesen, dass die Abstände von der Erde \emph{linear} bestimmt werden, sobald man das Verhältnis dieser drei Größen als bekannt ansieht. Es ist also ein Fehler, wenn Herr \textsc{Kramp} eine Aufgabe, deren \emph{strenge} Lösung nur eine Gleichung ersten Grades erfordert, auf eine kubische Gleichung zurückführt, und dies selbst unter Vernachlässigung einiger Glieder, bei deren Berücksichtigung das Verfahren von Herrn \textsc{Kramp} sogar eine biquadratische Gleichung ergeben hätte. Noch wichtiger ist aber der zweite Punkt. In Art. 131 des angeführten Werks ist bereits ausführlich bewiesen, dass es unzulässig ist, jene drei Größen $n''$, $n$, $n'$ \emph{zugleich} proportional zu den Zeitabständen $t$, $t'$, $t+t'$ zu setzen, da das Ergebnis nicht einmal als Näherung gelten darf.

Falls Herr \textsc{Kramp}, wie es fast scheint, die Theoria motus corporum coelestium nicht kannte, die fünf Jahre vor der Blockade von Straßburg erschienen ist, hätte er sich doch leicht davon überzeugen können, dass sein Verfahren im Kern auf einem Grundsatz beruht, dessen Unhaltbarkeit schon lange bekannt ist. Die drei oben erwähnten Größen sind nämlich, wie man bei einiger Aufmerksamkeit sofort sieht, lediglich die doppelten Flächen der Dreiecke zwischen dem ersten und zweiten, zweiten und dritten, ersten und dritten Radiusvektor. Setzt man also diese drei Größen \emph{zugleich} proportional zu den drei Zeitabständen, so nimmt man lediglich an, dass die drei Positionen des Kometen auf einer Geraden liegen, deren beide Abschnitte zu den Zeitabständen proportional sind. Ebenso ist die sogenannte équation essentielle von Herrn \textsc{Kramp} lediglich die Bedingungsgleichung dafür, dass die drei Positionen auf einer Geraden liegen. Es ist aber seit Langem bekannt, dass diese Voraussetzung beim Kometenproblem einerseits zu einer Lösung führen muss, die nur Gleichungen ersten Grades erfordert, und andererseits selbst als Näherung unzulässig ist. Einen ebenso gründlichen wie klaren Beweis findet man insbesondere in der vortrefflichen Schrift von \textsc{Olbers}, die Herr \textsc{Kramp} selbst in der vorliegenden Abhandlung anführt.
